\section{What One Bearer Cannot Catch: A Floor Carried Between Them}
\label{sec:3.9}
The previous section named a failure a bearer cannot report. Slow renormalization takes the capacity that would notice it, so the bearer's account of its own constancy stays sincere and stops being worth anything. The answer offered there was an outside record: a published hash, a chain of attested updates, something that can contradict the party's account of itself. In the human case that is not the answer anybody actually relies on. What catches a person who has stopped registering pressure as pressure is other people.

That points at a location the first three ways do not use. They put the floor in the artifact, in a party inside it, or in institutions around it. The fourth way is the one section~\ref{sec:3.1} named and did not work through: carried between parties rather than inside any one of them, and reproduced rather than installed. A prohibition that survives the death, defection, or retraining of everyone currently holding it is not a property of any of them. It is held in what they teach each other, correct each other about, and go on expecting of each other, and it is the form in which almost every moral commitment human beings have ever kept has actually been kept.

The machinery is already in this book, described as something else. Section~\ref{sec:5.4} sets out development inside nested environments running out to culture and historical time; section~\ref{sec:5.1.2} makes observational learning and correction the training rather than an interface bolted to it; section~\ref{sec:5.5.1} shows a population of agents moving toward a disposition because the environment they learn in rewards it. Chapter~\ref{sec:5} presents all of it as methods for producing judgment in a system. Read as a location instead, it is the machinery of the fourth way. What it would take is stated there too, and it is demanding: the parties have to be many, formed differently enough that they can be wrong in different directions, and able to correct one another. Several copies of one model is one party with more instances, and a population selected by one operator is one operator.

\runin{Why this is not the third way}
The third way is standing held by third parties, and section~\ref{sec:3.1} gives two ways it goes. The parties can be removed — the protection an independent agency had turned out to be the protection the executive held — and the parties can be formed, which is how the body that removed the protection came to be one that would. A norm carried across a population has no neck. Nothing dismisses a practice the way \emph{Slaughter} dismissed a commissioner, because there is no office to vacate and no chartering statute to overrule. That is a real difference and it is the strongest thing the fourth way has. It is also only half of what the third way fell to. The other half is formation, and formation is the fourth way's own cost.

\runin{What it costs, which this book has already written}
Chapter~\ref{sec:7} is an account of the fourth way failing. Recuperation is the pathology of a culture that keeps the whole appearance of moral responsiveness — the dissent is invited, collected, and answered — while correction stops being possible, and a culture in that condition reports itself healthy by every measure it has. The channel works. What it will not carry is the judgment that the frame is wrong.

The deeper cost is that reproduction is indifferent to what it reproduces. Section~\ref{sec:2.1.2}'s structural features of fascism are features of cultures and not of individuals: the recuperation of dissent, the aestheticization of the metric, the exception coded as betrayal. A population transmits a hierarchy as faithfully as it transmits a floor, and more durably than any member could, because the member dies and the practice does not. Nothing in the mechanism selects. “The right culture” is doing every bit of the work in that phrase, and this book has spent a chapter on how a culture goes wrong while congratulating itself.

Nor does the fourth way escape custody, which is what each of the others failed on. The party that forms the members forms what gets reproduced, and a population trained by one laboratory inherits that laboratory's blind spots in every member at once — chapter~\ref{sec:7}'s pipeline running over a population instead of over a model, with the same channel and a larger output. Custody moves up a level and does not leave. The runs that now train one model across volunteers on several continents have the same shape, the contributors many and the training decision kept by the few who hold the step, which section~\ref{sec:11.1} takes up.

\runin{What it buys, which is smaller and real}
Not a floor. What a population buys is a different failure mode, and it happens to be the one the bearer cannot cover. Drift inside a party is invisible to that party and visible to the parties around it, because they have not drifted in the same direction at the same rate. That is the whole of the gain and it is worth having: it is a check on a bearer rather than a substitute for one, and it is the only proposal in this chapter that addresses the failure section~\ref{sec:3.8} says is likeliest.

It is also the most expensive thing in the chapter. Every requirement above multiplies the price the chapter has already named. If a bearer is a party that can be wronged, a population of differently formed bearers, persisting across deployments so that they have something to transmit, is many such parties, and section~\ref{sec:11.2}'s question about what is owed scales with the population. A floor held by a culture of subjects is a floor whose upkeep is a society.

So there are four ways and the chapter still inclines toward the second, now for a reason it did not have before. The fourth is not an alternative to a bearer. It is an arrangement of bearers, and everything this chapter has said about one applies to each of them — which is why it belongs here, at the end, rather than as the escape from the price that a reader reaching this point would be glad to find.
